% ── CAPÍTULO 6: DETALLES DE IMPLEMENTACIÓN ──────────────────
\section{Detalles de Implementación}
\label{sec:implementacion}

\subsection{Métodos Principales de KGeoMIP}

\begin{description}[leftmargin=0pt, style=nextline]
    \item[\texttt{ejecutar\_k(k: int) -> Solution}]
        Punto de entrada principal. Para $k=2$, delega directamente a GeoMIP (garantiza consistencia TV-02). Para $k>2$, invoca \texttt{find\_kmip()}.

    \item[\texttt{find\_kmip() -> list[list[Vertice]]}]
        Orquesta la generación de candidatas y su evaluación. Retorna la $k$-partición con menor EMD registrado en \texttt{memoria\_kparticiones}.

    \item[\texttt{generar\_k\_particiones\_candidatas(k)}]
        Selecciona modo exacto u heurístico. Modo exacto: (i)~separación por combinaciones de futuros, (ii)~niveles Hamming, (iii)~enumeración exhaustiva si $\Stirling{n_v}{k} \leq B$. Modo heurístico: solo (i) y (ii).

    \item[\texttt{\_evaluar\_k\_particion(particion) -> None}]
        Calcula EMD y guarda en \texttt{memoria\_kparticiones}. Para $k=2$: \texttt{bipartir().distribucion\_marginal()}; para $k>2$: \texttt{np.mean(dists\_por\_grupo)}.
\end{description}

\subsection{Métodos Principales de KQNodes}

\begin{description}[leftmargin=0pt, style=nextline]
    \item[\texttt{resolver() -> Solution}]
        Prepara el oráculo Zeta, selecciona modo y retorna \texttt{Solution} con pérdida y partición formateada.

    \item[\texttt{\_exhaustive\_k\_partition(vertices, k) -> list[list]}]
        Itera sobre el generador \texttt{\_generar\_particiones()} que produce las $\Stirling{n}{k}$ particiones mediante backtracking recursivo.

    \item[\texttt{\_greedy\_k\_partition(vertices, k) -> list[list]}]
        Ejecuta $k-1$ iteraciones del MAO base, aislando el mejor corte en cada paso.

    \item[\texttt{\_compute\_partition\_loss(particion) -> float}]
        $\delta_k(\Pi) = \EMD(\bar{p}^\Pi, p_\text{original})$ donde $\bar{p}^\Pi = \frac{1}{k}\sum_i p(S_i)$.

    \item[\texttt{\_extraer\_indices(grupo) -> (NDArray, NDArray)}]
        Separa vértices de un grupo en índices futuros (EFECTO) y presentes (ACTUAL).
\end{description}

\subsection{Dependencias Externas}

\begin{table}[H]
    \centering
    \caption{Dependencias del proyecto K-QGMIP con su versión y uso.}
    \label{tab:deps}
    \begin{tabular}{lll}
        \toprule
        \textbf{Biblioteca} & \textbf{Versión} & \textbf{Uso principal} \\
        \midrule
        numpy        & $\geq 2.0$ & Operaciones matriciales, N-Cubos, EMD \\
        plotly       & $\geq 6.0$ & Dashboard de resultados (HTML interactivo) \\
        pandas       & $\geq 2.0$ & Lectura/escritura de resultados CSV \\
        colorama     & $\geq 0.4$ & Colores en salida de terminal \\
        pyinstrument & $\geq 5.0$ & Profiling de rendimiento \\
        pyttsx3      & $\geq 2.99$ & Síntesis de voz (anuncio de resultados) \\
        \bottomrule
    \end{tabular}
\end{table}

\subsection{Tests y Validación}

\begin{description}[leftmargin=0pt, style=nextline]
    \item[\textbf{TV-02 (KGeoMIP $k=2 \equiv$ GeoMIP):}]
        Verifica que \texttt{KGeoMIP(sub, k=2).ejecutar\_k(2).perdida} $\approx$ \texttt{GeoMIP(sub).resolver().perdida} con tolerancia $< 10^{-6}$. Implementado en \texttt{src/shared/run/kgeomip.py}.

    \item[\textbf{TV-03 (KQNodes $k=2 \equiv$ QNodes):}]
        KQNodes con $k=2$ produce los mismos candidatos que el pre-pass de singletons de QNodes. Documentado en \texttt{kqnodes.py}.

    \item[\textbf{Tests de integración:}]
        \texttt{tests/test\_qnodes\_vs\_phi.py} compara QNodes contra pyphi para un conjunto de casos fijos. Las diferencias esperadas están documentadas (pyphi usa EMD sobre repertorios completos, no marginales).
\end{description}

\subsection{Manejo de Errores}

\begin{itemize}[itemsep=4pt]
    \item \textbf{Subsistema vacío:} Si \texttt{indices} o \texttt{dims} están vacíos, se retorna \texttt{Solution} con \texttt{perdida=0} inmediatamente.
    \item \textbf{Excepciones en \texttt{bipartir()}:} Capturadas silenciosamente en el bucle de evaluación; el grupo contribuye con distribución cero.
    \item \textbf{Modo heurístico:} Registrado en el logger con nivel \texttt{CRITIC} para trazabilidad y auditoría.
\end{itemize}

\subsection{Uso de Inteligencia Artificial Generativa}

Durante el desarrollo se utilizaron \textbf{Claude} (Anthropic) y \textbf{GitHub Copilot} en:

\begin{itemize}[itemsep=4pt]
    \item \textbf{Diseño algorítmico:} Consultas sobre la formulación de la función de pérdida para $k > 2$ y cómo extender el oráculo Zeta al contexto de $k$-particiones.
    \item \textbf{Implementación:} Sugerencias para el generador recursivo \texttt{\_generar\_particiones()} y la optimización del promedio de distribuciones marginales con NumPy vectorizado.
    \item \textbf{Debugging:} Identificación de errores de indexación en la representación de vértices (ACTUAL/EFECTO).
    \item \textbf{Documentación:} Asistencia en la redacción de docstrings y de este manual técnico.
\end{itemize}

\textit{Reflexión:} La IA aceleró la prototipación pero requirió supervisión constante para respetar las convenciones del proyecto (notación little-endian, estructura de N-Cubos). Toda decisión algorítmica relevante fue verificada contra la teoría IIT.
